% =====================================================================
%  courses/trigonometria/semana_04_arcos/ejercicios.tex — examen y tarea: semana N° 02
% ---------------------------------------------------------------------
%  Cursos: ÁLGEBRA y TRIGONOMETRÍA
%  Compilar desde main.tex con LuaLaTeX
% =====================================================================

% --- ENCABEZADO DE PÁGINA — SEMANA 02 --------------------------------------
\newpage
\thispagestyle{fancy}

% --- CABECERA DECORATIVA ---------------------------------------------------
\begin{tikzpicture}[remember picture, overlay]
  \fill[AzulPizarra]
    ([xshift=0cm, yshift=0cm]current page.north west)
    rectangle ([xshift=0cm, yshift=-2.8cm]current page.north east);
  \fill[DoradoAcadémico]
    ([yshift=-2.8cm]current page.north west)
    rectangle ([yshift=-3.1cm]current page.north east);
  \node[anchor=west, text=white]
    at ([xshift=1.5cm, yshift=-1.0cm]current page.north west)
    {\large\sffamily\bfseries SEMANA N° 02 — EXAMEN Y TAREA};
  \node[anchor=west, text=GrisMedio]
    at ([xshift=1.5cm, yshift=-1.9cm]current page.north west)
    {\normalsize\sffamily\bfseries ÁLGEBRA \textcolor{DoradoAcadémico}{|} Polinomios y Funciones
     \quad\textcolor{DoradoAcadémico}{//}\quad
     \textbf{TRIGONOMETRÍA} \textcolor{DoradoAcadémico}{|} Área del Sector Circular};
  \node[anchor=east, text=white]
    at ([xshift=-1.5cm, yshift=-1.4cm]current page.north east)
    {\small\sffamily  2026};
\end{tikzpicture}

% --- BLOQUE A — ÁLGEBRA ----------------------------------------------------

% --- Separador de Curso ----------------------------------------------------
\begin{tcolorbox}[
  enhanced,
  colback=AzulPizarra, colframe=AzulPizarra,
  arc=3pt, boxrule=0pt,
  left=10pt, right=10pt, top=5pt, bottom=5pt
]
  \centering
  {\large\sffamily\bfseries\color{white} ÁLGEBRA}
  \quad\textcolor{DoradoAcadémico}{\rule{0.25\linewidth}{0.8pt}}
  \quad{\normalsize\sffamily\color{GrisMedio} Semana N° 02 · Polinomios y Funciones}
\end{tcolorbox}

\vspace{0.3cm}

% --- Bloque EXAMEN ÁLGEBRA -------------------------------------------------
\begin{tcolorbox}[
  enhanced,
  colback=AzulClaro!50, colframe=AzulMedio,
  arc=2pt, boxrule=0.7pt,
  fonttitle=\sffamily\bfseries\small\color{white},
  title={EXAMEN — ÁLGEBRA \hfill \small 5 preguntas · Tiempo: 20 min},
  attach boxed title to top left={yshift=-2mm, xshift=6mm},
  boxed title style={colback=AzulMedio, arc=2pt}
]
\small\sffamily
\textit{Instrucciones:} Marca la alternativa correcta. \textbf{La alternativa correcta aparece resaltada.}
\end{tcolorbox}

\vspace{0.2cm}

\begin{multicols}{2}
\setcounter{numejercicio}{0}

% --- Ejercicio E-A1 --------------------------------------------------------
\ejercicio Si $P(x) = 3x - 7$, calcule $P(3x^2 + 5)$.

\begin{alternativas}
  \item $25$
  \item \colorbox{AzulClaro!80}{$\mathbf{625}$}
  \item $600$
  \item $500$
  \item $340$
\end{alternativas}
\vspace{0.25cm}

% --- Ejercicio E-A2 --------------------------------------------------------
\ejercicio Calcular $P(P(x))$, si $P(x) = 2x + 1$. Hallar $P(P(1))$.

\begin{alternativas}
  \item $0$
  \item $64$
  \item \colorbox{AzulClaro!80}{$\mathbf{7}$}
  \item $4$
  \item $16$
\end{alternativas}
\vspace{0.25cm}

% --- Ejercicio E-A3 --------------------------------------------------------
\ejercicio Calcule $f(3)$, si:
\[
  2f(x) = x - 1 + f\!\left(\tfrac{x}{3}\right)
\]

\begin{alternativas}
  \item $1$
  \item \colorbox{AzulClaro!80}{$\mathbf{16}$}
  \item $4$
  \item $32$
  \item $8$
\end{alternativas}
\vspace{0.25cm}

% --- Ejercicio E-A4 --------------------------------------------------------
\ejercicio Calcule la suma de coeficientes de:
\[
  P(x) = (3x-2)^5 + (1-x)^n - (x-3)^2 + 7
\]

\begin{alternativas}
  \item $17$
  \item $48$
  \item $60$
  \item $20$
  \item \colorbox{AzulClaro!80}{$\mathbf{7}$}
\end{alternativas}
\vspace{0.25cm}

% --- Ejercicio E-A5 --------------------------------------------------------
\ejercicio Si $P(x+2) = x + P(x)$ y $P(3) = 1$. Halle $P(5) + P(1)$.

\begin{alternativas}
  \item $3$
  \item $2$
  \item \colorbox{AzulClaro!80}{$\mathbf{4}$}
  \item $0$
  \item $5$
\end{alternativas}

\end{multicols}

\vspace{0.4cm}

% --- Bloque TAREA ÁLGEBRA --------------------------------------------------
\begin{tcolorbox}[
  enhanced,
  colback=VerdeSuave!50, colframe=VerdeOliva,
  arc=2pt, boxrule=0.7pt,
  fonttitle=\sffamily\bfseries\small\color{white},
  title={TAREA — ÁLGEBRA \hfill \small 5 preguntas · Para casa},
  attach boxed title to top left={yshift=-2mm, xshift=6mm},
  boxed title style={colback=VerdeOliva, arc=2pt}
]
\small\sffamily
\textit{Instrucciones:} Resuelve en casa y presenta en la siguiente clase. Marca la alternativa correcta.
\end{tcolorbox}

\vspace{0.2cm}

\begin{multicols}{2}
\setcounter{numejercicio}{0}

% --- Tarea A1 --------------------------------------------------------------
\ejercicio Si $P(x) = ax^2 + 2x - 1$ y $P(-2) = 7$, halle $a$.

\begin{alternativas}
  \item $\tfrac{1}{2}$
  \item $\tfrac{2}{3}$
  \item $\tfrac{6}{5}$
  \item $0$
  \item $-1$
\end{alternativas}
\vspace{0.25cm}

% --- Tarea A2 --------------------------------------------------------------
\ejercicio Halle $a + b$, sabiendo que:
\[
  15 - 4x \equiv a(2-x) + b(1+x)
\]

\begin{alternativas}
  \item $0$
  \item $1$
  \item $\tfrac{1}{2}$
  \item $\tfrac{3}{2}$
  \item $2$
\end{alternativas}
\vspace{0.25cm}

% --- Tarea A3 --------------------------------------------------------------
\ejercicio Si $P(2x+3) = 7 - 6x$, halle $P(x+1)$.

\begin{alternativas}
  \item $16$
  \item $32$
  \item $30$
  \item $64$
  \item $-3x + 10$
\end{alternativas}
\vspace{0.25cm}

% --- Tarea A4 --------------------------------------------------------------
\ejercicio Hallar $a + b$, si son términos semejantes:
$9x^{2a+1}$ y $-2x^9 y^{5b-3}$

\begin{alternativas}
  \item $7$
  \item $10$
  \item $15$
  \item $4$
  \item $12$
\end{alternativas}
\vspace{0.25cm}

% --- Tarea A5 --------------------------------------------------------------
\ejercicio Sea $P(x) = x^2 - 3x + 2$. Calcule $P(2) \cdot P(0) + P(1)$.

\begin{alternativas}
  \item $3$
  \item $2$
  \item $1$
  \item $6$
  \item $0$
\end{alternativas}

\end{multicols}

\vspace{0.5cm}

% --- BLOQUE B — TRIGONOMETRÍA (Área del Sector Circular) -------------------

% --- Separador de Curso ----------------------------------------------------
\begin{tcolorbox}[
  enhanced,
  colback=AzulPizarra, colframe=AzulPizarra,
  arc=3pt, boxrule=0pt,
  left=10pt, right=10pt, top=5pt, bottom=5pt
]
  \centering
  {\large\sffamily\bfseries\color{white} TRIGONOMETRÍA}
  \quad\textcolor{DoradoAcadémico}{\rule{0.2\linewidth}{0.8pt}}
  \quad{\normalsize\sffamily\color{GrisMedio} Semana N° 02 · Área del Sector Circular}
\end{tcolorbox}

\vspace{0.3cm}

% --- Bloque EXAMEN TRIGONOMETRÍA -------------------------------------------
\begin{tcolorbox}[
  enhanced,
  colback=AzulClaro!50, colframe=AzulMedio,
  arc=2pt, boxrule=0.7pt,
  fonttitle=\sffamily\bfseries\small\color{white},
  title={EXAMEN — TRIGONOMETRÍA \hfill \small 5 preguntas · Tiempo: 20 min},
  attach boxed title to top left={yshift=-2mm, xshift=6mm},
  boxed title style={colback=AzulMedio, arc=2pt}
]
\small\sffamily
\textit{Instrucciones:} Marca la alternativa correcta. \textbf{La alternativa correcta aparece resaltada.}
\end{tcolorbox}

\vspace{0.2cm}

\begin{multicols}{2}
\setcounter{numejercicio}{0}

% --- Ejercicio E-T1 --------------------------------------------------------
\ejercicio En el sector circular, el ángulo central mide $\dfrac{\pi}{3}$ rad y el radio mide $6$ cm. Calcule el área.

\begin{center}
\begin{tikzpicture}[scale=0.7]
  \fill[AzulClaro!60] (0,0) -- (30:2.2) arc(30:90:2.2) -- cycle;
  \draw[AzulPizarra, thick] (0,0) circle(2.2);
  \draw[AzulPizarra, thick] (0,0) -- (30:2.2);
  \draw[AzulPizarra, thick] (0,0) -- (90:2.2);
  \draw[AzulMedio, ->] (0.5,0) arc(0:60:0.5) node[midway, right, font=\tiny] {$\frac{\pi}{3}$};
  \node[font=\tiny, AzulPizarra] at (55:1.4) {$S$};
  \node[font=\tiny] at (60:2.6) {$r=6$};
\end{tikzpicture}
\end{center}

\begin{alternativas}
  \item $3\pi\text{ cm}^2$
  \item \colorbox{AzulClaro!80}{$\mathbf{6\pi\text{ cm}^2}$}
  \item $9\pi\text{ cm}^2$
  \item $12\pi\text{ cm}^2$
  \item $4\pi\text{ cm}^2$
\end{alternativas}
\vspace{0.25cm}

% --- Ejercicio E-T2 --------------------------------------------------------
\ejercicio El arco de un sector circular mide $L = 4$ cm y el radio mide $R = 8$ cm. Calcule el área $S$.

\begin{center}
\begin{tikzpicture}[scale=0.75]
  \fill[AzulClaro!60] (0,0) -- (20:2.6) arc(20:70:2.6) -- cycle;
  \draw[AzulPizarra, thick] (0,0) -- (20:2.6) node[right, font=\tiny]{$R=8$};
  \draw[AzulPizarra, thick] (0,0) -- (70:2.6);
  \draw[AzulMedio, very thick] (20:2.6) arc(20:70:2.6) node[midway, right, font=\tiny]{$L=4$};
  \node[font=\tiny, AzulPizarra] at (45:1.4) {$S$};
\end{tikzpicture}
\end{center}

\begin{alternativas}
  \item $8\text{ cm}^2$
  \item $32\text{ cm}^2$
  \item $20\text{ cm}^2$
  \item \colorbox{AzulClaro!80}{$\mathbf{16\text{ cm}^2}$}
  \item $12\text{ cm}^2$
\end{alternativas}
\vspace{0.25cm}

% --- Ejercicio E-T3 --------------------------------------------------------
\ejercicio Calcula el área del sector circular con ángulo central $\theta = 1$ rad y radio $R = 8$ cm.

\begin{center}
\begin{tikzpicture}[scale=0.75]
  \fill[AzulClaro!60] (0,0) -- (10:2.4) arc(10:67.3:2.4) -- cycle;
  \draw[AzulPizarra, thick] (0,0) -- (10:2.4) node[right, font=\tiny]{$R=8$};
  \draw[AzulPizarra, thick] (0,0) -- (67.3:2.4);
  \draw[AzulMedio, ->] (0.6,0) arc(0:57.3:0.6) node[midway, above, font=\tiny]{$1$ rad};
  \node[font=\tiny, AzulPizarra] at (38:1.3) {$S$};
\end{tikzpicture}
\end{center}

\begin{alternativas}
  \item $16\text{ cm}^2$
  \item $64\text{ cm}^2$
  \item \colorbox{AzulClaro!80}{$\mathbf{32\text{ cm}^2}$}
  \item $8\text{ cm}^2$
  \item $48\text{ cm}^2$
\end{alternativas}
\vspace{0.25cm}

% --- Ejercicio E-T4 --------------------------------------------------------
\ejercicio Se tiene un sector circular de área $S$. Si el ángulo central se triplica y el radio se duplica, ¿cuál es el nuevo área?

\begin{center}
\begin{tikzpicture}[scale=0.78]
  % Sector original
  \fill[AzulClaro!40] (0,0) -- (15:1.5) arc(15:75:1.5) -- cycle;
  \draw[AzulPizarra] (0,0) -- (15:1.5);
  \draw[AzulPizarra] (0,0) -- (75:1.5);
  \draw[AzulMedio, ->] (0.45,0) arc(0:60:0.45) node[midway, right, font=\tiny]{$\alpha$};
  \node[font=\tiny] at (45:0.85) {$S$};
  \node[font=\tiny] at (45:1.75) {$m$};
  % Flecha
  \draw[->, NaranjaTierra, thick] (2.2,0.5) -- (3.2,0.5) node[midway, above, font=\tiny]{$\times 12$};
  % Sector nuevo
  \fill[NaranjaSuave!80] (4.5,0) -- ($(4.5,0)+(15:2.8)$) arc(15:75:2.8) -- cycle;
  \draw[NaranjaTierra] (4.5,0) -- ($(4.5,0)+(15:2.8)$) node[right, font=\tiny]{$2m$};
  \draw[NaranjaTierra] (4.5,0) -- ($(4.5,0)+(75:2.8)$);
  \draw[NaranjaTierra!80!black, ->] ($(4.5,0)+(0.55,0)$) arc(0:60:0.55)
    node[midway, right, font=\tiny]{$3\alpha$};
  \node[font=\tiny] at ($(4.5,0)+(45:1.6)$) {$A$};
\end{tikzpicture}
\end{center}

\begin{alternativas}
  \item $3S$
  \item $6S$
  \item $9S$
  \item \colorbox{AzulClaro!80}{$\mathbf{12S}$}
  \item $24S$
\end{alternativas}
\vspace{0.25cm}

% --- Ejercicio E-T5 --------------------------------------------------------
\ejercicio Calcula el área de la región sombreada. Los radios de los sectores son $OB = 6$ cm y $OD = 2\sqrt{3}$ cm, con ángulo $\angle BOC = 15°$.

\begin{center}
\begin{tikzpicture}[scale=0.85]
  % Sector grande OB
  \fill[AzulClaro!70]
    (0,0) -- ({cos(15)*2.4},{sin(15)*2.4}) arc(15:165:2.4) -- (0,0);
  % Quitar sector pequeño OD
  \fill[white]
    (0,0) -- ({cos(15)*1.38},{sin(15)*1.38}) arc(15:165:1.38) -- (0,0);
  % Sombreado final (corona)
  \fill[AzulClaro!70]
    ({cos(15)*2.4},{sin(15)*2.4}) arc(15:165:2.4) --
    ({cos(165)*1.38},{sin(165)*1.38}) arc(165:15:1.38) -- cycle;
  % Contornos
  \draw[AzulPizarra, thick] (0,0) -- ({cos(15)*2.4},{sin(15)*2.4})
    node[right, font=\tiny]{$B$};
  \draw[AzulPizarra, thick] (0,0) -- ({cos(165)*2.4},{sin(165)*2.4})
    node[left, font=\tiny]{$C$};
  \draw[AzulPizarra] (0,0) -- ({cos(15)*1.38},{sin(15)*1.38})
    node[right, font=\tiny]{$D$};
  \draw[AzulPizarra] (0,0) -- ({cos(165)*1.38},{sin(165)*1.38});
  \draw[AzulMedio, thick] ({cos(15)*2.4},{sin(15)*2.4}) arc(15:165:2.4);
  \draw[AzulMedio] ({cos(15)*1.38},{sin(15)*1.38}) arc(15:165:1.38);
  \node[font=\tiny] at (0.0, -0.2) {$O$};
  \node[font=\tiny, AzulPizarra] at (90:1.85) {$S$};
  % Ángulo
  \draw[AzulMedio,->] (0.55,0) arc(0:15:0.55) node[right, font=\tiny]{$15°$};
  % Cotas
  \node[font=\tiny] at ({cos(15)*1.9},{sin(15)*1.9}) {$6$};
  \node[font=\tiny] at ({cos(15)*0.69+0.05},{sin(15)*0.69}) {$2\!\sqrt{3}$};
\end{tikzpicture}
\end{center}

\begin{alternativas}
  \item $2\pi\text{ cm}^2$
  \item $3\pi\text{ cm}^2$
  \item $4\pi\text{ cm}^2$
  \item \colorbox{AzulClaro!80}{$\mathbf{\pi\text{ cm}^2}$}
  \item $\tfrac{\pi}{2}\text{ cm}^2$
\end{alternativas}

\end{multicols}

\vspace{0.4cm}

% --- Bloque TAREA TRIGONOMETRÍA --------------------------------------------
\begin{tcolorbox}[
  enhanced,
  colback=VerdeSuave!50, colframe=VerdeOliva,
  arc=2pt, boxrule=0.7pt,
  fonttitle=\sffamily\bfseries\small\color{white},
  title={TAREA — TRIGONOMETRÍA \hfill \small 5 preguntas · Para casa},
  attach boxed title to top left={yshift=-2mm, xshift=6mm},
  boxed title style={colback=VerdeOliva, arc=2pt}
]
\small\sffamily
\textit{Instrucciones:} Resuelve en casa y presenta en la siguiente clase. Marca la alternativa correcta.
\end{tcolorbox}

\vspace{0.2cm}

\begin{multicols}{2}
\setcounter{numejercicio}{0}

% --- Tarea T1 --------------------------------------------------------------
\ejercicio En un sector circular, el ángulo central mide $18°$ y el radio mide $20$ cm. ¿Cuál es su área?

\begin{center}
\begin{tikzpicture}[scale=0.78]
  \fill[VerdeSuave!80] (0,0) -- (0:2.2) arc(0:18:2.2) -- cycle;
  \draw[VerdeOliva, thick] (0,0) -- (0:2.2) node[right, font=\tiny]{$20$ cm};
  \draw[VerdeOliva, thick] (0,0) -- (18:2.2);
  \draw[AzulMedio, ->] (0.6,0) arc(0:18:0.6) node[midway, above right, font=\tiny]{$18°$};
  \node[font=\tiny, VerdeOliva] at (9:1.2) {$S$};
\end{tikzpicture}
\end{center}

\begin{alternativas}
  \item $\dfrac{10\pi}{9}\text{ cm}^2$
  \item $20\pi\text{ cm}^2$
  \item $\dfrac{100\pi}{9}\text{ cm}^2$
  \item $10\pi\text{ cm}^2$
  \item $\dfrac{200\pi}{9}\text{ cm}^2$
\end{alternativas}
\vspace{0.25cm}

% --- Tarea T2 --------------------------------------------------------------
\ejercicio El arco de un sector circular mide $L = 3$ cm y el radio $R = 10$ cm. Halle el área $S$.

\begin{center}
\begin{tikzpicture}[scale=0.78]
  \fill[VerdeSuave!80] (0,0) -- (5:2.3) arc(5:40:2.3) -- cycle;
  \draw[VerdeOliva, thick] (0,0) -- (5:2.3) node[right, font=\tiny]{$R=10$};
  \draw[VerdeOliva, thick] (0,0) -- (40:2.3);
  \draw[VerdeOliva!80!black, very thick] (5:2.3) arc(5:40:2.3)
    node[midway, right, font=\tiny]{$L=3$};
  \node[font=\tiny, VerdeOliva] at (22:1.3) {$S$};
\end{tikzpicture}
\end{center}

\begin{alternativas}
  \item $10\text{ cm}^2$
  \item $\dfrac{3}{2}\text{ cm}^2$
  \item $30\text{ cm}^2$
  \item $15\text{ cm}^2$
  \item $\dfrac{15}{2}\text{ cm}^2$
\end{alternativas}
\vspace{0.25cm}

% --- Tarea T3 --------------------------------------------------------------
\ejercicio Si AOD y COB son sectores circulares concéntricos (ver figura), con $OA = 3$, $OC = 5$, y ángulo $\theta$. Calcule $k = \dfrac{S_1}{S_2}$.

\begin{center}
\begin{tikzpicture}[scale=0.80]
  % Sector grande
  \fill[AzulClaro!40] (0,0) -- (10:2.5) arc(10:70:2.5) -- cycle;
  % Sector pequeño (blanco encima)
  \fill[white] (0,0) -- (10:1.5) arc(10:70:1.5) -- cycle;
  % Corona sombreada = S2
  \fill[AzulClaro!70] (10:1.5) arc(10:70:1.5) -- (70:2.5) arc(70:10:2.5) -- cycle;
  % Triángulo S1
  \fill[VerdeSuave!80] (0,0) -- (10:1.5) arc(10:70:1.5) -- cycle;
  % Contornos
  \draw[AzulPizarra] (0,0) -- (10:2.5) node[right, font=\tiny]{$B$};
  \draw[AzulPizarra] (0,0) -- (70:2.5) node[above, font=\tiny]{$C$};
  \draw[AzulMedio] (0,0) -- (10:1.5) node[below, font=\tiny]{$A$};
  \draw[AzulMedio] (0,0) -- (70:1.5) node[left, font=\tiny]{$D$};
  \draw[AzulMedio] (10:1.5) arc(10:70:1.5);
  \draw[AzulPizarra, thick] (10:2.5) arc(10:70:2.5);
  \node[font=\tiny, AzulPizarra] at (40:0.9) {$S_1$};
  \node[font=\tiny] at (40:2.1) {$S_2$};
  \node[font=\tiny] at (-0.15,-0.15) {$O$};
  % Cotas
  \node[font=\tiny] at ($(0,0)!0.5!(10:1.5)+(0,-0.15)$) {$3$};
  \node[font=\tiny] at ($(10:1.5)!0.5!(10:2.5)+(0,-0.12)$) {$2$};
\end{tikzpicture}
\end{center}

\begin{alternativas}
  \item $k = \dfrac{9}{16}$
  \item $k = \dfrac{3}{5}$
  \item $k = \dfrac{9}{25}$
  \item $k = \dfrac{3}{4}$
  \item $k = \dfrac{1}{3}$
\end{alternativas}
\vspace{0.25cm}

% --- Tarea T4 --------------------------------------------------------------
\ejercicio Se tiene un sector circular de área $S$. Si el ángulo central se duplica y el radio se triplica, ¿cuál es el área del nuevo sector circular?

\begin{center}
\begin{tikzpicture}[scale=0.80]
  \fill[VerdeSuave!60] (0,0) -- (10:1.3) arc(10:60:1.3) -- cycle;
  \draw[VerdeOliva] (0,0) -- (10:1.3) node[right, font=\tiny]{$m$};
  \draw[VerdeOliva] (0,0) -- (60:1.3);
  \draw[VerdeOliva!80, ->] (0.4,0) arc(0:50:0.4)
    node[midway, right, font=\tiny]{$\alpha$};
  \node[font=\tiny] at (35:0.7) {$S$};
  % Flecha
  \draw[->, NaranjaTierra, thick] (1.8,0.4) -- (2.7,0.4);
  % Nuevo
  \fill[NaranjaSuave!80] (3.8,0) -- ($(3.8,0)+(10:2.4)$) arc(10:60:2.4) -- cycle;
  \draw[NaranjaTierra] (3.8,0) -- ($(3.8,0)+(10:2.4)$) node[right, font=\tiny]{$3m$};
  \draw[NaranjaTierra] (3.8,0) -- ($(3.8,0)+(60:2.4)$);
  \draw[NaranjaTierra!80, ->] ($(3.8,0)+(0.55,0)$) arc(0:50:0.55)
    node[midway, right, font=\tiny]{$2\alpha$};
  \node[font=\tiny] at ($(3.8,0)+(35:1.4)$) {$A$};
\end{tikzpicture}
\end{center}

\begin{alternativas}
  \item $6S$
  \item $9S$
  \item $12S$
  \item $18S$
  \item $24S$
\end{alternativas}
\vspace{0.25cm}

% --- Tarea T5 --------------------------------------------------------------
\ejercicio En la figura, AOB y COD son sectores circulares. Si $OC = 6$ cm, $L_{COD} = 6$ y el sector AOB tiene $OA = x$ y $S_{AOB} = 4 \cdot S_{COD}$, calcule $x$.

\begin{center}
\begin{tikzpicture}[scale=0.85]
  % Sector COD (pequeño)
  \fill[VerdeSuave!80] (0,0) -- (5:1.6) arc(5:65:1.6) -- cycle;
  \draw[VerdeOliva] (0,0) -- (5:1.6) node[right, font=\tiny]{$C$};
  \draw[VerdeOliva] (0,0) -- (65:1.6) node[above, font=\tiny]{$D$};
  \draw[VerdeOliva, thick] (5:1.6) arc(5:65:1.6);
  \node[font=\tiny, VerdeOliva] at (35:0.9) {$S_2$};
  \node[font=\tiny] at (0.08,1.7) {$L=6$};
  % Sector AOB (grande)
  \fill[AzulClaro!40] (0,0) -- (5:2.8) arc(5:65:2.8) -- cycle;
  \fill[white] (0,0) -- (5:1.6) arc(5:65:1.6) -- cycle; % hueco
  \draw[AzulPizarra, thick] (0,0) -- (5:2.8) node[right, font=\tiny]{$B$};
  \draw[AzulPizarra, thick] (0,0) -- (65:2.8) node[above, font=\tiny]{$A$};
  \draw[AzulPizarra, thick] (5:2.8) arc(5:65:2.8);
  \node[font=\tiny, AzulPizarra] at (35:2.2) {$S_1$};
  \node[font=\tiny] at (-0.15,-0.15) {$O$};
  \node[font=\tiny] at ($(0,0)!0.6!(5:1.6)+(0,-0.15)$) {$6$};
  \node[font=\tiny] at ($(0,0)!0.75!(5:2.8)+(0,-0.15)$) {$x$};
\end{tikzpicture}
\end{center}

\begin{alternativas}
  \item $x = 6$
  \item $x = 10$
  \item $x = 8$
  \item $x = 12$
  \item $x = 9$
\end{alternativas}

\end{multicols}

\vspace{0.5cm}

% --- CLAVE DE RESPUESTAS ---------------------------------------------------
\begin{tcolorbox}[
  enhanced, colback=GrisClaro, colframe=GrisMedio,
  arc=2pt, boxrule=0.5pt,
  fonttitle=\sffamily\bfseries\small\color{GrisCarbón},
  title={Clave de Respuestas — Exámenes (Álgebra y Trigonometría)}
]
\small\sffamily
\begin{multicols}{2}
  \textbf{Álgebra:}\quad
  1.~B \quad 2.~C \quad 3.~B \quad 4.~E \quad 5.~C \\[4pt]
  \textbf{Trigonometría:}\quad
  1.~B \quad 2.~D \quad 3.~C \quad 4.~D \quad 5.~D
\end{multicols}
\end{tcolorbox}

% FIN DEL ARCHIVO ejercicios_examen_tarea.tex
